\section{The actual exceptional coefficient primes and the ramified CM endpoint}
\label{sec:cubic-exceptional-coefficients}

The exceptional primes in a source hypothesis and the exceptional
summands in a numerical all-prime sum are different notions.  This
section makes the distinction for the fixed cubic character of
$269$.  The source paths below have the pinned root
\path{Development/AutomorphicArithmetic/AutomorphicVarieties/}
and commit \texttt{7cbf7a0c32fca1badf05931879a37c3f8a2b03b4}.

\begin{proposition}[AC-K29-exceptional: exact source exclusion and fixed-three scope]
\label{AC-K29-exceptional}
For $K=\mathbb Q(\omega)$, the literal hypothesis
\texttt{HeightAndVanishing/Counts.lean}, \texttt{GD}, excludes
coefficient primes $2$ and $3$.  It does not exclude $269$:
$\ell=269$, $e=2$, $N=2$, and the trivial Dirichlet character
satisfy every clause of \texttt{GD}.
This assertion concerns the datum, not the existence of a Type II
point of its deformation ring.

For the fixed strict groups of AC-G26, with coefficient
$(\mathbb Q_3(\omega)/\mathbb Z_3[\omega])
 (\chi^{\pm1}\varepsilon_3^{-e})$, AC-K23 gives the exact
unnormalized length
\[
 m_e=4+1_{e\equiv2\pmod3},\qquad
 \sum_{e=1}^{n}m_e\log3
   =\left(4n+\left\lfloor\frac{n+1}{3}\right\rfloor\right)\log3.
                                                               \tag{CE.1}
\]
In the normalized convention (GK.8a), each character contributes
$m_e/2$, and the two cubic characters together contribute $m_e$.
This closes that fixed character family at three.  Formula (CE.1)
does not identify the groups with an additional nontrivial
Dirichlet twist, or the groups on a growing host field.
\end{proposition}
\begin{proof}
The fields of \texttt{GD} are printed explicitly in
\texttt{Counts.lean}.  Its \texttt{l\_odd} is $\ell\ne2$.
Its \texttt{l\_cop} requires
$\ell\nmid |\operatorname{disc}K|h_K$; here
$|\operatorname{disc}K|=3$ and $h_K=1$, so $\ell=3$ is
excluded.  Independently its \texttt{gen} is
$\ell-1\nmid2e$, which also fails at three for every $e$.

For the stated example at $269$, \texttt{l\_odd},
\texttt{e\_pos}, \texttt{N\_even}, \texttt{quad},
\texttt{imag}, \texttt{l\_cop} and \texttt{cop} hold directly.
The character has order one, so \texttt{ord} holds; its
$269$-conductor component is trivial, so \texttt{ine} holds.
The exact definition
\texttt{DirichletPrimeToLPartMulLPowerPart.lean},
\texttt{PsiEven}, is $\psi(-1)=(-1)^e$, which holds for $e=2$.
The same file defines \texttt{chiP}; it is trivial here.
For \texttt{eta}, inertia at three acts trivially on the
mod-$269$ cyclotomic character and on \texttt{chiP}, but
nontrivially on the ramified quadratic extension
$K/\mathbb Q$.  Thus their joint kernel is not contained
in $G_K$, exactly as \texttt{eta} requires.
Finally $268\nmid4$, proving \texttt{gen}.

The positive cubic character case of (CE.1) is AC-K23.
Conjugation of the coefficient ring exchanges $\chi$ and
$\chi^{-1}$, fixes the cyclotomic values, and transports all
the same global and local cochain kernels.  The two lengths
are therefore equal.  Since
$\operatorname{length}_{\mathbb Z_3[\omega]}
 (\mathbb Z_3[\omega]/3)=2$, the normalized length is
$m_e/2$.  The local condition and field in this argument
are precisely those in AC-G26.  An additional character
changes the coefficient and usually its splitting field;
no comparison with that different group is supplied by
(CE.1).
\end{proof}

\begin{proposition}[AC-K29-269: the ramified cubic Artin representation has no crystalline refinement]
\label{AC-K29-269}
Put $p=269$ and
$\rho=\operatorname{Ind}_{G_K}^{G_{\mathbb Q}}\chi$.
At $p$, after adjoining the cubic values to the coefficient
field, one has
\[
 \rho|_{I_p}=\xi\oplus\xi^{-1},\qquad
 \operatorname{ord}(\xi)=3,\qquad
 \rho^{I_p}=0,\qquad a_p(\rho)=2,
 \qquad D_{\rm cris,\mathbb Q_p}(\rho)=0.               \tag{CE.2}
\]
Here $a_p(\rho)$ is the Artin conductor exponent.
The primitive theta series $g=\theta_\chi$ has local
Euler factor $L_p(g,s)=1$ and $U_p$ eigenvalue zero.
For every finite-order character $\tau$ of
$G_{\mathbb Q_p}$ and every integer $n$,
\[
        D_{\rm cris,\mathbb Q_p}(\rho\tau(n))=0.        \tag{CE.3}
\]
Thus the ordinary or finite-slope CM quotient used in
AC-K29-inert and AC-K29-irregular is unavailable at this
particular classical theta point.  This conclusion is
about those constructions, not a nonexistence theorem
for some other Euler system.
\end{proposition}
\begin{proof}
The extension $K_p/\mathbb Q_p$ is unramified quadratic,
since $269\equiv-1\pmod3$.  The polynomial
$X^3-269$ is Eisenstein over $K_p$ and $\mu_3\subset K_p$.
It defines a totally tamely ramified cyclic extension of
degree three.  Consequently $\chi$ on inertia is a
nontrivial character of exact order three.  Induction
restricts to the two characters $\xi,\xi^p=\xi^{-1}$.
They are different and neither is trivial.  Frobenius
exchanges their lines, so the local induced representation
is irreducible.  Wild inertia acts trivially.  These facts
prove the first four assertions of (CE.2), including
$a_p(\rho)=2-\dim\rho^{I_p}=2$.

A finite-image representation is potentially unramified
and potentially crystalline with all Hodge--Tate weights
zero.  To see its crystalline invariants here explicitly,
pass to a finite Galois extension $L/\mathbb Q_p$ killing
its finite image.  Then
$D_{\rm cris,L}(\rho)=L_0\otimes\rho$, where $L_0$ is
the maximal unramified subfield.  Inertia in
$\operatorname{Gal}(L/\mathbb Q_p)$ acts trivially on
$L_0$ and has no fixed vector on $\rho$.  Taking invariants
therefore gives $D_{\rm cris,\mathbb Q_p}(\rho)=0$.
This is the usual finite-representation case of crystalline
Galois descent, using $B_{\rm cris}^{G_L}=L_0$.
The exact period-ring input is Brinon--Conrad,
\emph{CMI summer school notes on $p$-adic Hodge theory},
\url{https://math.stanford.edu/~conrad/papers/notes.pdf},
Theorem 9.1.5 and the immediately preceding invariant-field
consequence, printed p.~131.  Example 6.3.9, printed
pp.~81--82, supplies the finite-inertia de Rham statement.
These standard period-ring theorems are cited inputs;
the finite-group descent calculation above is the proof
needed for this particular representation.

The Artin local factor is the determinant on inertia
invariants, hence is one.  The primitive theta form has
that same local factor: its Dirichlet-series construction
omits the unique prime ideal above $269$, which divides
the conductor of $\chi$.  In particular every coefficient
$a_{269m}(g)$ is zero.  The Fourier-series definition of
$U_p$ therefore gives $U_pg=0$ directly, without a
numerical slope approximation.

For a rational local finite-order character $\tau$, local
class field theory identifies its inertia restriction with
a finite character of $\mathbb Z_p^\times$.  Its tame
part has order dividing $p-1=268$, and its wild part has
$p$-power order.  Neither can cancel a nontrivial character
of order three.  Thus both characters in
$(\rho\tau)|_{I_p}$ remain nontrivial; the finite-image
argument just given proves (CE.3) when $n=0$.
Tate twisting multiplies the crystalline period module by
a power of Fontaine's period $t$, so it preserves its
dimension, proving (CE.3) for every $n$.

The exact finite-slope source is Loeffler--Zerbes,
\emph{Rankin--Eisenstein classes in Coleman families},
\url{https://arxiv.org/abs/1506.06703}, published version
DOI \texttt{10.1186/s40687-016-0077-6}.
Remark 4.5.2, printed p.~27, requires $U_p$ invertible on
a finite-slope part.  Definition 4.6.1, printed p.~28,
requires $a_p$ invertible in the coefficient affinoid.
The weight-one extension, Theorem 4.7.2 and its complete
proof, printed pp.~29--30, is for a $p$-regular
stabilization; Remark 7.1.6, printed p.~42, expressly
notes that this weight-one specialization is ordinary.
The complete construction in Theorem 5.4.2, printed
pp.~36--37, normalizes its level-$p^n$ classes by
$(a_p(\mathcal F)a_p(\mathcal G))^{-n}$.
The zero $U_p$ eigenvalue above satisfies none of these
refinement conditions.  Raising the $p$-power level does
not repair this eigenform: all its old translates still
have only the zero generalized $U_p$ eigenvalue, since
$U_p(g(p^r z))=g(p^{r-1}z)$ for $r\ge1$ and $U_pg=0$.
Thus changing the number of Coleman maps within that
finite-slope construction does not supply its missing
coefficient quotient.
\end{proof}

\paragraph{The precise remaining arithmetic input at $269$.}
Let $W=V|_{G_K}\chi^a$ be the actual mixed coefficient and
$U=W^*(1)$.  Once the corresponding actual local point
construction has been supplied, the same rank-two
Euler-system and duality argument would suffice if one
constructs ray norm-compatible classes for $U$, with one
fixed lattice denominator, together with two rational
base classes whose local images span
$H_s^1(K_p,U)$.  The norm polynomial must be the actual
$\det(1-X\operatorname{Fr}_{\mathfrak q}^{-1}\mid W)$,
and the classes must be Bloch--Kato away from $p$.
The rank-two image hypotheses and the strict dual bound
must remain those proved in the preceding Euler-system
modules.  Under these inputs, local Poitou--Tate duality
makes $H_f^1(K,W)$ strict at $p$, and the nonzero ray
Euler class kills that strict dual group, exactly as in
AC-K29-lower.  This is a sufficient missing input, not
an assertion that such classes have been constructed.
Two local linear regulator maps applied to a single
class do not by themselves give two independent global
classes or prove this spanning condition.

\paragraph{Numerical and realization scope.}
Deleting finitely many primes from (GK.9) cannot bound
their omitted summands.  Formula (CE.1) supplies the
fixed-three summand only for its exact character and
field; no estimate here controls arbitrary twists at
$2$, $3$, or $269$ on a growing host.  In particular no
bound on Iwasawa $\mu$ or $\lambda$ alone is being
converted into a bound at an arithmetic specialization.
The local representation in (CE.2) is the actual Artin
realization of the finite Kummer cover.  Its absent
inertia invariants and its absent crystalline periods
are distinct, explicitly computed properties; neither
is a claim about the spectral or prolate Hilbert spaces.
